\section*{Chapitre \ref{ch:hf:extensions-of-the-inventory-framework} --- Extensions du cadre d'inventaire}

\begin{solution}{exo:hf:extensions-of-the-inventory-framework:1}
$c=\sqrt{0.5e/210}=0.0804$~; $\delta^b(2)=0.667+2.5c=0.868$, $\delta^a(2)=0.667-1.5c=0.546$.
\end{solution}

\begin{solution}{exo:hf:extensions-of-the-inventory-framework:2}
$q_0=\mu/(2\phi)=1$ unité.
\end{solution}

\begin{solution}{exo:hf:extensions-of-the-inventory-framework:3}
$17^3=4\,913$ points pour trois instruments~; $17^{10}\approx2.0\times10^{12}$ pour dix~: la grille est hors d'atteinte, d'où les approximations quadratiques.
\end{solution}

\begin{solution}{exo:hf:extensions-of-the-inventory-framework:4}
$\delta^a+\delta^b=2/k+c\bigl((2q+1)-(2q-1)\bigr)/2=2/k+c$, indépendant de $q$~: $1.333+0.080=1.414$.
\end{solution}

\begin{solution}{exo:hf:extensions-of-the-inventory-framework:5}
Avec une corrélation positive, une position longue sur l'actif~1 compense une partie de la variance de la position courte sur l'actif~2, de sorte qu'acheter
l'actif~1 réduit la pénalité $\phi q^\top\Sigma q$. Avec une corrélation négative, la position compensatrice est une position courte sur l'actif~1~: le teneur de marché
affiche alors un cours acheteur moins élevé et un cours vendeur plus agressif.
\end{solution}

\begin{solution}{exo:hf:extensions-of-the-inventory-framework:6}
Le mouvement réévalue l'ensemble de la position restante, et non seulement l'unité échangée. Après une exécution à l'achat, la position vaut $q+1$, après une exécution à la vente $q-1$~; en moyenne
sur les positions que détient le teneur de marché, la perte par exécution, au-delà de l'unité elle-même, se réduit à un terme quadratique en $q$ que la fonction
valeur absorbe, et il ne reste que la moitié du mouvement par exécution.
\end{solution}

\begin{solution}{exo:hf:extensions-of-the-inventory-framework:7}
En $q=0$~: cours acheteur 0.543, cours vendeur 0.872 (le teneur de marché cherche à acheter)~; en $q=2$~: 0.708 des deux côtés, exactement les profondeurs sans dérive en $q=0$. La dérive déplace
l'ensemble de la grille de $q_0=2$.
\end{solution}

\begin{solution}{exo:hf:extensions-of-the-inventory-framework:8}
Trois erreurs sont possibles. Élargir de la totalité du mouvement au lieu de sa moitié (un mouvement permanent réévalue aussi la position)~; un pas de temps assez grossier pour modifier la réponse
(le pas grossier du chapitre ne favorisait aucun élargissement, le pas fin en favorisait la moitié du mouvement)~; et évaluer sur le P\&L moyen alors que la politique maximise le P\&L
diminué de la pénalité. Élargir de la moitié du mouvement permanent, vérifier avec un quart du pas de temps, et évaluer l'objectif.
\end{solution}

\begin{solution}{pb:hf:extensions-of-the-inventory-framework:1}
\begin{enumerate}
\item Une limite stricte $|q|\le\bar q$~; elle rend l'espace des états fini, de sorte que le système linéaire admet une plus grande valeur propre qui domine lorsque l'horizon croît.
\item Voir la \cref{prop:hf:extensions-of-the-inventory-framework:stationary} (Perron--Frobenius et la décomposition spectrale de $e^{M(T-t)}$).
\item Voir la \cref{prop:hf:extensions-of-the-inventory-framework:closed}.
\item Lorsque $\beta=\sqrt{\phi ke/A}$ est petit~; les écarts maximaux pour $|q|\le10$ valent 0.034, 0.70 et 5.0.
\item Un seul vecteur d'inventaire et la charge $\phi q^\top\Sigma q$ sur l'ensemble du livre.
\item La substitution $h=\tfrac1k\ln\omega$ linéarise le terme de chaque instrument seulement si tous partagent le même $k$.
\item 0.572, 0.700 et 0.828.
\item 1\,000 trajectoires communes de cinq unités de temps, milieux corrélés à 0.8, six pénalités~; le risque d'inventaire est la moyenne temporelle de $q^\top\Sigma q$.
\item 21~\% de moins (24.3 contre 30.8 pour un profit moyen de 673.3)~; 10~\% à 22~\% selon les pénalités.
\item La variance du P\&L provient aussi du caractère aléatoire des exécutions et de la capture du spread, que la cotation conjointe ne modifie pas.
\item La dérive intervient dans le problème de contrôle, par l'intermédiaire de la valeur de la détention d'inventaire.
\item Voir la \cref{prop:hf:extensions-of-the-inventory-framework:target}.
\item 355.6 contre 335.7 (5.9~\% de plus), avec une détention moyenne de 1.97 unité pour une cible de 2.
\item Moins~: le teneur de marché ne peut détenir que jusqu'à sa borne, et la récompense $-\phi(q-q_0)^2$ est alors maximale à la borne.
\item Voir la \cref{prop:hf:extensions-of-the-inventory-framework:adverse}.
\item Profit 262.7, 267.2, 258.9~; objectif 253.1, 258.6, 251.2 pour aucun, 0.15 et 0.3.
\item Avec une cotation fixe exécutée avec une probabilité d'environ 0.24 à chaque pas, la simulation discrète s'écarte du modèle en temps continu que résout la proposition~;
avec un quart du pas, elle s'accorde avec lui.
\item Sélection adverse~: elle est présente sur tous les marchés réels (le teneur de marché du chapitre~1 a perdu son spread à cause d'elle), et l'ajustement est une constante par côté.
\item Mesurer le mark-out de ses propres exécutions à l'horizon où la courbe se stabilise (livre~7, chapitre~23)~: le mark-out stabilisé vaut $\xi$.
\item Chacune constitue une modification d'une seule matrice~: sa taille (bornes, actifs), sa diagonale (pénalités, dérive) ou ses termes hors diagonale (sélection adverse).
\end{enumerate}
\end{solution}

\begin{solution}{iq:hf:extensions-of-the-inventory-framework:1}
Détenir le fonds, c'est être long de l'indice~: il faut décaler les cotations du contrat à terme comme si l'on en était déjà long, plus enclin à vendre et moins enclin à acheter, par
le ratio de couverture multiplié par l'inventaire du fonds.
\iqlookfor{une réflexion transversale au livre, avec un ratio de couverture.}
\end{solution}

\begin{solution}{iq:hf:extensions-of-the-inventory-framework:2}
La solution est $e^{M(T-t)}z$~; loin de l'horizon, la plus grande valeur propre domine, et les profondeurs ne dépendent que des rapports du vecteur propre dominant.
\iqlookfor{l'argument spectral, ou une intuition de contrôle ergodique.}
\end{solution}

\begin{solution}{iq:hf:extensions-of-the-inventory-framework:3}
Le mouvement est-il un coût sur l'unité (élargir de 0.3) ou un déplacement permanent du mid (élargir de 0.15)~? Il faut mesurer la courbe de mark-out et l'endroit où elle se stabilise,
et vérifier quelle part du mouvement revient.
\iqlookfor{la distinction et un plan de mesure.}
\end{solution}

\begin{solution}{iq:hf:extensions-of-the-inventory-framework:4}
Précalculer~: résoudre hors ligne, stocker les profondeurs dans une table indexée par l'inventaire (et par quelques classes de signal), ou utiliser la forme fermée~; rafraîchir la table
lorsque les paramètres changent, hors du hot path.
\iqlookfor{le précalcul et une consultation sur le hot path.}
\end{solution}

\begin{solution}{iq:hf:extensions-of-the-inventory-framework:5}
La corrélation entre les positions~: des positions qui se compensent paraissent risquées et des positions concentrées paraissent diversifiées. Mesurer $q^\top\Sigma q$ (ou une
version à modèle factoriel) sur l'ensemble du livre.
\iqlookfor{la variance de portefeuille, et non une somme de variances.}
\end{solution}

\begin{solution}{iq:hf:extensions-of-the-inventory-framework:6}
La récompense instantanée $\mu q-\phi q^2=-\phi(q-\mu/(2\phi))^2+\mu^2/(4\phi)$~: le problème est celui à dérive nulle dans la variable décalée $q-q_0$, dont
l'optimum centre l'inventaire en zéro.
\iqlookfor{compléter le carré.}
\end{solution}
